\documentclass{article}
\usepackage[T1]{fontenc}
\usepackage{amsmath}
\usepackage{amssymb}

\title{Assignment 1: Math1014}

\begin{document}
\setlength{\parindent}{0pt}
\maketitle

\section{Linear Algebra}

\subsection{Question 1}
we are given the vector space
\[
V_N = \operatorname{span} \left\{\cos(2\pi kt),\sin(2\pi kt)\mid 0\leq k\leq N\right\}.
\]

\subsubsection*{(a)}
two functions are linearly independent if the only scalars $a,b$ satisfying
\[
a\cos(2\pi t) + b \sin(2\pi t) = 0 \quad \text{for all } t \in [0,1]
\]
are $a = b = 0$. if we can find even one nontrivial choice of $t$ that forces $a=0$, and another that forces $b=0$, we're done.

- at $t=0$,
\[
a\cos(0) + b\sin(0) = a(1) + b(0) = a = 0
\]
so $a$ must be $0$
 
- at $t = \frac{1}{4}$
\[
a\cos\left(\tfrac{\pi}{2}\right) + b\sin\left(\tfrac{\pi}{2}\right) = a(0) + b(1) = b = 0
\]
so $b$ must be $0$ too
 
since the equation forces $a=b=0$, the only linear combination equal to the zero function is the trivial one. 
hence $\cos(2\pi t)$ and $\sin(2\pi t)$ are \textbf{linearly independent} in $V_1$.

\subsubsection*{(b)}
we need the dimension of $V_1$

the definition says
\[
V_N= \operatorname{span} \left\{ \cos(2\pi kt),\sin(2\pi kt) \mid0\leq k\leq N \right\}
\]

if $N=1$, then $k$ can be either $0$ or $1$.

so we get

\[
V_1= \operatorname{span} \left\{ \cos(0),\sin(0), \cos(2\pi t),\sin(2\pi t) \right\}
\]

now, $\cos(0)=1$ and $\sin(0)=0$

so the span becomes:

\[
V_1= \operatorname{span} \left\{ 1,\cos(2\pi t),\sin(2\pi t) \right\}
\]

the zero function doesn't add anything to the span, so we can ignore it

we therefore have three functions:
\[
1,\cos(2\pi t),\sin(2\pi t)
\]

these are linearly independent, so they form a basis for \(V_1\).

therefore, \boxed{dim(V_1) = 3}

\subsubsection*{(c)}
finding the dimension of $V_2$,

when $N=2$, the possible values of $k$ are $0,1,2$

therefore,
\[
V_2 = \operatorname{span} \left\{\cos(0), \sin(0), cos(2\pi t), sin(2\pi t), cos(4\pi t), sin(4\pi t) \right\}
\]

again, $\cos(0) = 1$ and $\sin(0) = 0$, so (leaving out the $0$ terms),
\[
V_2 = \operatorname{span} \left\{1, cos(2\pi t), sin(2\pi t), cos(4\pi t), sin(4\pi t) \right\}
\]

there are now 5 independent functions. therefore,
\[
\boxed{dim(V_2) = 5}
\]

for a general $N$, we always have:
\begin{itemize}
    \item one constant function $1$,
    \item $N$ cosine functions corresponding to $k=1,...,N$,
    \item $n$ sine functions coreresponding to $K=1,...,N$
\end{itemize}
so, the total number of basis functions is
\[
1+N+N
\]
therefore,
\[
\boxed{dim(V_N) = 2N + 1}
\]

\subsection{Question 2}

we have $S:V_N \rightarrow R_m$ defined by
\[
S(f) = 
\begin{bmatrix}
    f(t_1) \\
    f(t_2) \\
    \vdots \\
    f(t_m)
\end{bmatrix}
\]

\subsubsection*{(a)}
we are given $m=3$ with sampling points $t_1=0$, $t_2=0.5$, $t_3=1$ and 
$$f(t) = 3+\cos(2\pi t)$$

since $S$ records the value of $f$ at each sampling point,
\[
S(f)=
\begin{bmatrix}
    f(0) \\
    f(0.5) \\
    f(1)
\end{bmatrix}
\]
\begin{itemize}
    \item at $t=0$, $$f(0) = 3+\cos(2\pi(0))$$ $$f(0)=3+\cos(0)$$
    since $\cos(0)=1$, $$f(0) = 4$$
    \item at $t=0.5$, $$f(0.5) = 3+\cos(2\pi(0.5))$$ $$f(0.5)=3+\cos(\pi)$$
    since $\cos(\pi)=-1$, $$f(0.5) = 3-1=2$$
    \item at $t=1$, $$f(1) = 3+\cos(2\pi(1))$$ $$f(1)=3+\cos(2\pi)$$
    since $\cos(2\pi)=1$, $$f(1) = 3=1=4$$
\end{itemize}

putting all these together into the matrix, we get
\[
\boxed{
    S(f) = 
    \begin{bmatrix}
        4\\
        2\\
        4
    \end{bmatrix}
}
\]

\subsubsection*{(b)}
now we have $$g(t) = \sin(12\pi t)$$
again,
\[
S(g) = 
\begin{bmatrix}
    g(0)\\
    g(0.5)\\
    g(1)
\end{bmatrix}
\]

\begin{itemize}
    \item at $t=0$, $$g(0)=\sin(12\pi(0))=\sin(0) = 0$$
    \item at $t=05$, $$g(0.5)=\sin(12\pi(0.5))=\sin(6\pi) = 0$$
    \item at $t=1$, $$g(1)=\sin(12\pi(1))=\sin(12\pi) = 0$$
\end{itemize}

putting all these together into the matrix, we get
\[
\boxed{
    S(f) = 
    \begin{bmatrix}
        0\\
        0\\
        0
    \end{bmatrix}
}
\]

\subsubsection*{(c)}
for a function to be a linear map, we need
\[
S(af+bg)=aS(f)+bS(g) \quad \text{for any functions } f, g \in V_N \text{ and scalars } a, b \in \mathbb{R}
\]
starting with the LHS,
\[
S(f) = 
\begin{bmatrix}
    (af+bg)(t_1)\\
    (af+bg)(t_2)\\
    \vdots\\
    (af+bg)(t_m)
\end{bmatrix}
\]
we know that at each sample point $t_i$, the addition and scalar multiplication is
\[
(af+bg)(t_i)=af(t_i)+bg(t_i)
\]
therefore,
\[
S(f) = 
\begin{bmatrix}
    af(t_1) + bg(t_1)\\
    af(t_2) + bg(t_2)\\
    \vdots\\
    af(t_m) + bg(t_m)
\end{bmatrix}
\]
separating into two vectors:
\[
S(f) = 
a
\begin{bmatrix}
    f(t_1)\\
    f(t_2)\\
    \vdots\\
    f(t_m)
\end{bmatrix}
+
b
\begin{bmatrix}
    g(t_1)\\
    g(t_2)\\
    \vdots\\
    g(t_m)
\end{bmatrix}
\]
we can see that these two vectors are exactly $S(f)$ and $S(g)$

therefore, $$S(af+bg) = aS(f)+bS(g)$$
hence, $$\boxed{\text{S is linear}}$$

\subsubsection*{(d)}
an element of $V_N$ can be represented as a vector in $\mathbb{R}^{2N+1}$

from question 1, $dim(V_N)=2N+1$, so $[f]_\mathcal{B} \in \mathbb{R}^{2N+1}$.

so there are $2N+1$ functions in this basis, meaning $[f]_\mathcal{B}$ has $2N+1$ components

$S$ takes $f$ and evaluates it at $m$ sample points

\[
S(f) =
\begin{bmatrix}
    f(t_1)\\
    f(t_2)\\
    \vdots\\
    f(t_m)
\end{bmatrix}
\]

therefore, $S(f) \in \mathbb{R}^m$. so the output has $m$ components

so we need a matrix $A$ such that

\[
\text{matrix }A \times \text{input vector} = \text{output vector}
\]

therefore, $A$ needs to take a vector with $2N+1$ entries and turn it into a vector with $m$ entries.

this means $A$ needs $2N+1$ columns (one for each basis function) and $m$ rows (one for each sample point).

therefore,

\[
\boxed{A\in\mathbb{R}^{m\times(2N+1)}}
\]

\subsubsection*{(e)}
we know from 2(d) that

\[
A\in\mathbb{R}^{m\times(2N+1)}
\]

so $A$ has $m$ rows and $2N+1$ columns.

the rank of a matrix cannot be bigger than the number of rows or the number of columns.

therefore,

\[
rank(A)\leq m
\]

and

\[
rank(A)\leq 2N+1
\]

so the upper bound is the smaller of these two numbers:

\[
\boxed{rank(A)\leq\min(m,2N+1)}
\]

this is because the rank tells us how many independent rows or columns there are, so we cannot have more independent rows than $m$, or more independent columns than $2N+1$.

\subsubsection*{(f)}
we can use the rank-nullity theorem:
\[
rank(A)+nullity(A)=\text{number of columns of }A
\]

from 2(d), $A$ has $2N+1$ columns, so
\[
rank(A)+nullity(A)=2N+1
\]

rearranging,
\[
nullity(A)=2N+1-rank(A)
\]

from part (e),
\[
rank(A)\leq\min(m,2N+1)
\]

so the smallest possible nullity occurs when the rank is as large as possible.
therefore,
\[
nullity(A)\geq 2N+1-\min(m,2N+1)
\]

so,
\[
\boxed{nullity(A)\geq\max(0,2N+1-m)}
\]

if $m<2N+1$, then
\[
\boxed{nullity(A)\geq2N+1-m>0}
\]

\subsubsection*{(g)}
from part(f) we know that
\[
nullity(A)\geq2N+1-m.
\]

we want
\[
dim(null(A))>0.
\]

this will definitely happen when
\[
2N+1-m>0.
\]

so,
\[
2N+1>m.
\]

therefore,
\[
\boxed{m<2N+1}
\]

leads to
\[
\boxed{dim(null(A))>0}.
\]

in other words, if we have fewer sample points than the number of basis functions, there must be some non-zero functions that get mapped to the zero vector.

\subsubsection*{(h)}

\section{Calculus}

\subsection{Question 3}

\subsubsection*{a}

\subsubsection*{b}

\subsubsection*{c}

\subsection{Question 4}

\subsubsection*{a}

\subsubsection*{b}

\subsubsection*{c}

\end{document}